    %%%%%%%%%%%%%%%%%%%%%%%%%%%%%%%%%%%%%%%%%%%%%%%%%%%%%%%%%%
  %% MARK: art: The Uniqueness of the Quotient and the Physical Choice
  %%%%%%%%%%%%%%%%%%%%%%%%%%%%%%%%%%%%%%%%%%%%%%%%%%%%%%%%%%
  
  \begin{article}
    \artlabel[The Uniqueness of the Quotient and the Physical Choice]
    \label{art: The Physical Choice}
    
    The framework of diffeology forces a re-evaluation of the relationship between physical modeling and geometric structure.
    Consider the case of a composite system,
    such as the Deuteron,
    modeled as two indistinguishable spins.
    The classical phase space for distinguishable spins is $M = S^2 \times S^2$.
    Imposing indistinguishability leads to the quotient $\Sigma_4 = M / \fS_2$.
    
    In a context lacking the tools to handle singularities,
    one might observe that $\Sigma_4$ is homeomorphic to the complex projective plane $P^2(\CC)$.
    Since $P^2(\CC)$ carries a standard smooth manifold structure,
    it is tempting to assume that nature ``chooses'' this smooth structure by default,
    effectively smoothing out the singularity of the quotient.
    This leads to the standard quantum mechanical description of the deuteron as a single elementary particle (spin 1).
    
    However,
    in diffeology,
    this ambiguity disappears.
    The quotient $\Sigma_4$,
    equipped with the quotient diffeology,
    is a well-defined,
    rigid geometric object.
    It is distinct from $P^2(\CC)$ despite their topological equivalence;
    they are not diffeomorphic,
    and as discussed regarding the ``Boman Paradox'' \cite{PIZ25},
    topology and algebra are insufficient to define geometry.
    
    Therefore,
    we cannot transform $\Sigma_4$ into $P^2(\CC)$ by mere will or for analytical convenience.
    The choice is not between two geometries for the same system,
    but between two distinct physical hypotheses:
    \begin{enumerate}
      \item \textbf{The Elementary Hypothesis:} The system is a single,
      indivisible entity.
      The space is the homogeneous $P^2(\CC)$.
      The quantization yields an irreducible representation.
      \item \textbf{The Composite Hypothesis:} The system is composed of two identical subsystems.
      The space is the inhomogeneous $\Sigma_4$.
      The quantization yields a structure that retains the memory of the diagonal singularity (the state of coincidence).
    \end{enumerate}
    Diffeology ensures that the geometry remains faithful to the physical hypothesis.
    If one chooses the composite model,
    one must accept the singular geometry of $\Sigma_4$ and its resulting prequantum groupoid.
    The ``smoothing'' of the singularity is not a geometric regularization but a modification of the physical model itself.
  \end{article}
  
  %%%%%%%%%%%%%%%%%%%%%%%%%%%%%%%%%%%%%%%%%%%%%%%%%%%%%%%%%%
  %% MARK: art: Example Loops S3
  %%%%%%%%%%%%%%%%%%%%%%%%%%%%%%%%%%%%%%%%%%%%%%%%%%%%%%%%%%
  
  \begin{article}
    \artlabel[Example: The Prequantum Groupoid for $\pmb{\Loops(\S^3)}$]
    \label{art: Example Loops S3}
    
    The construction of the prequantum groupoid $(\TT_\omega, \blambda)$ is applicable not only to singular spaces but also to infinite-dimensional diffeological spaces,
    provided they satisfy the conditions of the main theorem (connected, simply connected, with a closed $2$-form having discrete periods).
    A relevant example is the loop space of the three-sphere,
    $\X = \Loops(\S^3)$.
    
    The space $\X = \Loops(\S^3)$ is an infinite-dimensional diffeological space.
    Its path-connectedness is given by $\pi_0(\Loops(\S^3)) \simeq \pi_1(\S^3) = \{0\}$;
    thus $\X$ is connected.
    Its simple connectedness is given by $\pi_1(\Loops(\S^3)) \simeq \pi_2(\S^3) = \{0\}$;
    thus $\X$ is simply connected.
    
    Consider the volume $3$-form $\vol$ on $\S^3$.
    Let $\omega = \CHO\vol \in \Omega^2(\X)$ be the $2$-form on $\X$ obtained by applying the chain-homotopy operator.
    The $2$-form $\omega$ is closed,
    which follows from $d\vol = 0$ and the properties of the chain-homotopy operator.
    The group of periods $\P_\omega$ of the form $\omega$ is the group of periods of the $1$-form $\Komega \restriction \Loops(\Loops(\X))$.
    This group is a homomorphic image of $\pi_1(\Loops(\X)) \simeq \pi_2(\X) \simeq \pi_3(\S^3) = \ZZ$.
    The periods of $\omega$ are related to integrals of $\vol$ over cycles in $\pi_3(\S^3)$.
    For a suitable normalization of the volume form $\vol$ on $\S^3$ (e.g., such that $\int_{\S^3} \vol = 1$),
    the group of periods $\P_\omega$ is isomorphic to $\ZZ$.
    
    The torus of periods is $\T_\omega = \RR/\P_\omega = \RR/\ZZ \simeq \S^1$.
    
    According to the main theorem,
    the prequantum groupoid $\TT_\omega$ for $(\X, \omega)$ has $\X = \Loops(\S^3)$ as its objects.
    The space of morphisms is $\cY = \Paths(\X) / \sim_\omega = \Paths(\Loops(\S^3)) / \sim_\omega$,
    where the equivalence relation $\sim_\omega$ is determined by the $2$-form $\omega$ and its periods $\P_\omega = \ZZ$.
    
    The isotropy group at any object $\ell \in \X$ (i.e., at any loop $\ell \in \Loops(\S^3)$) is isomorphic to the torus of periods:
    $\TT_{\omega, \ell} \simeq \T_\omega \simeq \S^1$.
    This means that at each point in the infinite-dimensional space $\Loops(\S^3)$,
    the quantum phase information is captured by an $\S^1$ group.
    
    The prequantum $1$-form $\blambda$ on $\cY$ is the unique form such that $\class_\omega^*(\blambda) = \Komega$.
    
    Explicitly representing the space of morphisms $\cY = \Paths(\Loops(\S^3)) / \sim_\omega$ in a simpler form (like a product space) is generally not possible due to the infinite-dimensional nature of $\X$ and the non-exactness of $\omega$ (since $\P_\omega = \ZZ \neq \{0\}$).
    The space $\cY$ is genuinely the quotient space defined by the equivalence relation.
    
    This example illustrates how the prequantum groupoid construction extends to infinite-dimensional settings.
    The situation shares some analogy with the prequantization of the $2$-sphere $(\S^2, \omega_{S^2})$,
    where $\S^2$ is connected and simply connected,
    and the periods of the standard symplectic form are related to $\pi_2(\S^2) \simeq \ZZ$,
    leading to an $\S^1$ isotropy group.
    The key difference here is the infinite dimensionality of the base space $\X = \Loops(\S^3)$.
    
  \end{article}
  
  %%%%%%%%%%%%%%%%%%%%%%%%%%%%%%%%%%%%%%%%%%%%%%%%%%%%%%%%%%
  %% MARK: art: The Moment Map on the Prequantum Groupoid
  %%%%%%%%%%%%%%%%%%%%%%%%%%%%%%%%%%%%%%%%%%%%%%%%%%%%%%%%%%
  
  \begin{article}\artlabel[The Moment Map on the Prequantum Groupoid]
    \label{art:Momentum Map}
    
    Now that we have this symmetry group $\Aut(\TT_\omega, \blambda)$,
    we can define a moment map using its action,
    and here are its properties and what it tells us.
    
    The prequantum groupoid $(\TT_\omega, \blambda)$ provides a geometric framework to study the system $(\X, \omega)$.
    As established in Article \ref{art:Automorphisms Induced by Symmetries},
    its group of automorphisms $\G_\omega = \Aut(\TT_\omega, \blambda)$ is diffeologically isomorphic to the group of $\omega$-preserving diffeomorphisms of $\X$, $\Diff(\X, \omega)$.
    This group $\G_\omega$ acts smoothly on the space of morphisms $\cY = \Mor(\TT_\omega)$,
    preserving the prequantum $1$-form $\blambda$.
    
    According to the general theory of moment maps in diffeology developed in \cite{PIZ10} (see also \cite[Chap.\ 9]{PIZ13}),
    the existence of a $\G_\omega$-invariant $1$-form $\blambda$ on $\cY$ ensures the existence of a moment map $\Psi_\omega : \cY \to \cG_\omega^*$,
    where $\cG_\omega^*$ is the space of left-invariant $1$-forms on the diffeological group $\G_\omega$ (the space of momenta \cite[Art.~2.6]{PIZ10}).
    This map is defined for $y \in \cY$ by $\Psi_\omega(y) = \hat y^*(\blambda)$,
    where $\hat y : \G_\omega \to \cY$ is the orbit map $\hat y(\Phi) = \Phi(y)$.
    
    This moment map on the space of morphisms $\cY$ is directly related to the paths moment map defined in \cite[Art.~3.1]{PIZ10} for the action of $\G_\omega$ on $\X$.
    Let $\Psi_{\text{paths}} : \Paths(\X) \to \cG_\omega^*$ be the paths moment map given by $\Psi_{\text{paths}}(\gamma) = \hat\gamma^*(\Komega)$,
    where $\hat\gamma(\phi) = \phi \circ \gamma$.
    The map $\Psi_\omega$ is precisely the projection of $\Psi_{\text{paths}}$ onto the quotient $\cY$:
    
    \begin{proposition}{}
      The moment map $\Psi_\omega : \cY \to \cG_\omega^*$ satisfies $\Psi_\omega \circ \class_\omega = \Psi_{\text{paths}}$.
    \end{proposition}
    
    \begin{proof}
      Let $\gamma \in \Paths(\X)$,
      and let $y = \class_\omega(\gamma) \in \cY$.
      By definition of $\Psi_\omega$, $\Psi_\omega(y) = \hat y^*(\blambda)$,
      where $\hat y : \G_\omega \to \cY$ is the orbit map $\hat y(\Phi) = \Phi(y)$ for $\Phi \in \G_\omega$.
      Recalling that $\G_\omega = \Aut(\TT_\omega, \blambda)$ is isomorphic to $\Diff(\X,\omega)$ via $\Phi \mapsto \phi$ such that $\Phi([\gamma']_\omega) = [\phi \circ \gamma']_\omega$,
      the orbit map $\hat y$ is given by $\hat y(\Phi) = [\phi \circ \gamma]_\omega = \class_\omega(\phi \circ \gamma)$.
      Thus,
      $\hat y = \class_\omega \circ \hat\gamma$,
      where $\hat\gamma : \G_\omega \to \Paths(\X)$ is the orbit map $\hat\gamma(\phi) = \phi \circ \gamma$ as defined for $\Psi_{\text{paths}}$.
      Now we compute the pullback:
      $$
        \Psi_\omega(\class_\omega(\gamma)) = \hat y^*(\blambda) = (\class_\omega \circ \hat\gamma)^*(\blambda) = \hat\gamma^*(\class_\omega^*(\blambda)).
        $$
      By definition of $\blambda$ (Theorem \ref{sec:Main Theorem: The Prequantum Groupoid}, part 3),
      $\class_\omega^*(\blambda) = \Komega$.
      Substituting this gives:
      $$
        \Psi_\omega(\class_\omega(\gamma)) = \hat\gamma^*(\Komega).
        $$
      The right side is precisely $\Psi_{\text{paths}}(\gamma)$.
      Therefore, $\Psi_\omega \circ \class_\omega = \Psi_{\text{paths}}$.
    \end{proof}
    \smallskip
    
    The paths moment map $\Psi_{\text{paths}}$ is an additive homomorphism from the path concatenation to $(\cG_\omega^*, +)$ \cite[Art.~4.3]{PIZ10}.
    Since $\Psi_\omega$ factors through $\class_\omega$ and $\class_\omega$ preserves concatenation,
    $\Psi_\omega$ is an additive homomorphism from the groupoid $\TT_\omega$ to $(\cG_\omega^*, +)$:
    $$
      \Psi_\omega(y \cdot y') = \Psi_\omega(y) + \Psi_\omega(y')
      $$
    for any composable $y, y' \in \cY$.
    
    The additive property $\Psi_\omega(y \cdot y') = \Psi_\omega(y) + \Psi_\omega(y')$ implies that the value of $\Psi_\omega(y)$ for $y \in \cY$ depends only on the endpoints of $y$.
    This means $\Psi_\omega$ descends to a map $\psi_\omega : \X \times \X \to \cG_\omega^*$.
    This descent occurs if and only if $\Psi_\omega$ is constant on the equivalence classes $[\gamma]_\omega$ corresponding to loops,
    which is equivalent to $\Psi_\omega(\ell)=0$ for all loops $\ell \in \Loops(\X)$ (viewed as elements of $\cY$ via $\class_\omega$).
    The set of values $\Psi_{\text{paths}}(\ell) = \Psi_\omega(\class_\omega(\ell))$ for $\ell \in \Loops(\X)$ is the universal holonomy group $\Gamma_\omega = \Psi_{\text{paths}}(\Loops(\X))$ for the action of $\G_\omega$ on $\X$ \cite[Art.~3.7]{PIZ10}.
    
    Since $\X$ is connected and simply connected,
    its fundamental group $\pi_1(\X)$ is trivial.
    The universal holonomy group $\Gamma_\omega$ is an additive subgroup of $\cG_\omega^*$ which is a homomorphic image of $\pi_0(\Loops(\X))$ \cite[Art.~3.7]{PIZ10}.
    Since $\X$ is simply connected and connected,
    $\pi_0(\Loops(\X))$ is trivial.
    Therefore,
    for a simply connected $\X$,
    the universal holonomy $\Gamma_\omega$ is trivial,
    $\Gamma_\omega = \{0\}$.
    This property will not be preserved in the general case of $\pi_1(\X) \neq \{0\}$.
    In the general case,
    the universal holonomy $\Gamma_\omega$ will typically be non-trivial.
    
    The fact that $\Gamma_\omega=\{0\}$ for simply connected $\X$ implies that $\Psi_\omega(\ell)=0$ for all loops $\ell \in \Loops(\X)$,
    and thus $\Psi_\omega : \cY \to \cG_\omega^*$ descends to a well-defined map on $\X \times \X$.
    This descended map is the \emph{two-points moment map}:
    $$
      \psi_\omega : \X \times \X \to \cG_\omega^*
      \quad \text{defined by} \quad
      \psi_\omega(x,x') = \Psi_\omega(y),
      $$
    for any $y \in \cY$ with $\ends(y) = (x,x')$.
    This $\psi_\omega$ is the universal two-points moment map associated with $\omega$ (compare \cite[Art.~4.1]{PIZ10}).
    It satisfies the Chasles cocycle relation $\psi_\omega(x,x'') = \psi_\omega(x,x') + \psi_\omega(x',x'')$ for $x,x',x'' \in \X$.%
    \footnote{~This moment map $\psi_\omega$ can be interpreted as the moment map for the diffeological space $\X \times \X$ equipped with the $2$-form $\omega \ominus \omega = \pr_1^*\omega - \pr_2^*\omega$ and the diagonal action of $\G_\omega$.}
    
    From the two-points moment map $\psi_\omega$,
    we can derive the \emph{one-point moment map} $\mu_\omega : \X \to \cG_\omega^*$ by choosing a base point $x_0 \in \X$.
    As shown in \cite[Art.~5.1]{PIZ10},
    $\mu_\omega$ is a primitive of $\psi_\omega$,
    unique up to an additive constant $\varepsilon \in \cG_\omega^*$,
    satisfying $\psi_\omega(x,x') = \mu_\omega(x') - \mu_\omega(x)$.
    A standard choice is $\mu_\omega(x) = \psi_\omega(x_0,x)$.
    
    The equivariance properties of these moment maps and the associated Souriau's cocycle $\theta_\omega \in \Cinfty(\G_\omega, \cG_\omega^*)$ are described in detail in \cite[Art.~5.2]{PIZ10}.
    For simply connected $\X$,
    the holonomy $\Gamma_\omega$ is trivial,
    which means the target space for Souriau's cocycle and moment maps is $\cG_\omega^*$ (as opposed to $\cG_\omega^*/\Gamma_\omega$ in the general case).
    The universal Souriau class $\sigma_\omega$,
    which lives in $\H^1(\G_\omega, \cG_\omega^*/\Gamma_\omega)$,
    is therefore an element of $\H^1(\G_\omega, \cG_\omega^*)$.
    The triviality of this cohomology class implies the existence of an equivariant moment map $\mu_\omega$ \emph{op.~cit.}
    However,
    the triviality of $\Gamma_\omega$ does not in general imply the triviality of the entire cohomology group $\H^1(\G_\omega, \cG_\omega^*)$.
    Thus,
    for a simply connected $\X$,
    while the holonomy is trivial,
    Souriau's class $\sigma_\omega$ may be non-trivial,
    corresponding to non-trivial central extensions in the classical picture.
    
  \end{article}
  
  %%%%%%%%%%%%%%%%%%%%%%%%%%%%%%%%%%%%%%%%%%%%%%%%%%%%%%%%%%
  %% MARK: art: The Geometric Nature of Quantum States
  %%%%%%%%%%%%%%%%%%%%%%%%%%%%%%%%%%%%%%%%%%%%%%%%%%%%%%%%%%
  
  \begin{article}\artlabel[The Geometric Nature of Quantum States]
    \label{art:The Geometric Nature of Quantum States}
    
    The construction of the prequantum groupoid $\TT_\omega$ suggests a shift in the definition of a quantum state.
    Rather than starting with a Hilbert space of sections,
    we look for objects that are intrinsic to the groupoid structure itself.
    
    A natural candidate arises from the algebraic structure of the morphisms.
    Let $\V$ be a complex vector space (typically $\CC$).
    A smooth function $\Psi: \cY \to \V$ is called \emph{multiplicative} if,
    for any composable pair $y_1, y_2 \in \cY$,
    we have:
    $$
      \Psi(y_1 \cdot y_2) = \Psi(y_1) \cdot \Psi(y_2).
      $$
    This definition directly encodes the quantum phase.
    For any loop $\tau$ in the isotropy group $\TT_{\omega, x} \simeq \T_\omega$,
    the condition implies $\Psi(\tau \cdot y) = \chi(\tau)\Psi(y)$,
    where $\chi$ is a character of the torus of periods.
    
    These multiplicative functions are the geometric "atoms" of the quantum theory.
    They correspond to the "pure states" or "semiclassical amplitudes" associated with the geometry.
    While the full analytical machinery (superposition,
    operators) requires the construction of a convolution algebra (as discussed in Part II),
    these functions constitute the fundamental geometric dual of the prequantum groupoid.
    They demonstrate that the quantum state is not an external addition to the theory,
    but a covariant representation of the space of motions itself.
  \end{article}
  
  %%%%%%%%%%%%%%%%%%%%%%%%%%%%%%%%%%%%%%%%%%%%%%%%%%%%%%%%%%
  %% MARK: art: From Geometry to Algebra
  %%%%%%%%%%%%%%%%%%%%%%%%%%%%%%%%%%%%%%%%%%%%%%%%%%%%%%%%%%
  \begin{article}\artlabel[From Geometry to Algebra]
    \label{art:From Geometry to Algebra}
    
    The prequantum groupoid $\TT_\omega$ constitutes the objective geometric reality of the quantum system.
    However,
    to connect this geometry with the analytical machinery of quantum mechanics (operators,
    spectra,
    states),
    we must transition from the \emph{points} of the space to the \emph{functions} defined upon it.
    This transition occurs in two stages:
    the identification of elementary amplitudes and the construction of the algebra of observables.
    
    1.~\emph{Elementary Amplitudes (Multiplicative Functions).}
    A natural class of objects arises directly from the groupoid structure.
    A smooth,
    complex-valued function $\Psi$ on the space of morphisms $\cY$ is called \emph{multiplicative} if,
    for any composable pair $y, y' \in \cY$,
    we have:
    $$
      \Psi(y \cdot y') = \Psi(y) \cdot \Psi(y').
      $$
    This condition enforces that $\Psi$ transforms according to a character of the isotropy group $\T_\omega$.
    Specifically,
    for any loop $\tau \in \TT_{\omega, \src(y)} \simeq \T_\omega$,
    we have $\Psi(\tau \cdot y) = \chi(\tau)\Psi(y)$.
    These functions represent the "pure phases" or "semiclassical states" of the system.
    They are the geometric "atoms" of the theory,
    corresponding to flat sections in the bundle picture.
    
    2.~\emph{The Prequantum Algebra.}
    Physical states and observables,
    however,
    rely on the principle of linear superposition.
    They are not single paths,
    but weighted sums (integrals) over the space of motions.
    We define the \emph{Prequantum Algebra} $\cA_\omega$ as the space of compactly supported complex densities on $\cY$.
    The groupoid structure induces a natural \emph{convolution product} on this space:
    $$
      (f * g)(y) = \int_{y_1 \cdot y_2 = y} f(y_1) g(y_2).
      $$
    This operation transforms the geometry of the groupoid into a non-commutative $*$-algebra.
    The multiplicative functions defined above play a distinguished role here:
    they are the elementary characters (or "plane waves") of this algebra.
    The general elements of the algebra are linear combinations of these geometric phases.
    
    3.~\emph{The Threshold of Analysis.}
    The algebra $\cA_\omega$ retains the full covariance of the classical system.
    The group of symmetries $\Diff(X, \omega)$ acts by automorphisms on the groupoid,
    and therefore by automorphisms on the algebra.
    
    The transition to standard quantum mechanics requires a final step:
    \emph{Polarization}.
    This consists of selecting a specific representation of this algebra on a Hilbert space.
    This step is analytically necessary but geometrically traumatic,
    as it often breaks the symmetry of the groupoid.
    By stopping at the algebra $\cA_\omega$,
    we preserve the purity of the geometric signal.
    We have constructed the universal container for the physics,
    standing exactly at the threshold where geometry yields to analysis.
  \end{article}
  
  %%%%%%%%%%%%%%%%%%%%%%%%%%%%%%%%%%%%%%%%%%%%%%%%%%%%%%%%%%
  %% MARK: art: Relation to Classical Central Extensions
  %%%%%%%%%%%%%%%%%%%%%%%%%%%%%%%%%%%%%%%%%%%%%%%%%%%%%%%%%%
  
  \begin{article}
    \artlabel[Relation to Classical Central Extensions]
    \label{art:Automorphism Group Extension} % Keep the label for potential internal references
    
    It is natural to question how the isomorphism $\Aut(\TT_\omega, \blambda) \simeq \Diff(\X, \omega)$ relates to the well-known phenomenon of non-trivial central extensions arising in the classical geometric quantization of symmetries, even on simply connected symplectic manifolds.
    For instance,
    the prequantization of the standard symplectic structure on $\RR^{2n}$ leads to the Heisenberg group extension for the translation group,
    and the prequantization of the sphere $\S^2$ leads to the spin group $\mathrm{Spin}(3) \simeq \mathrm{SU}(2)$ as a central extension of $\SO(3)$.
    
    The key to understanding this lies in distinguishing the object to which the symmetries are lifted:
    
    $\ast$~In the \textbf{classical prequantum bundle picture},
    central extensions emerge when considering the lift of a symmetry group (or its Lie algebra) acting on the base manifold $\X$ to the total space of the \emph{prequantum principal $\U(1)$-bundle} $P \to \X$,
    or equivalently, to the space of sections of the associated line bundle. The failure of the lifted Lie bracket to match the lift of the base Lie bracket results in a cocycle valued in the Lie algebra of the fiber group ($\mathfrak{u}(1)$),
    leading to a central extension of the symmetry Lie algebra and group by $\U(1)$.
    
    $\ast$~In the \textbf{prequantum groupoid picture},
    the isomorphism $\Aut(\TT_\omega, \blambda) \simeq \Diff(\X, \omega)$ states that the full group of $\omega$-preserving diffeomorphisms $\phi$ lifts \emph{isomorphically} to act as automorphisms $\Phi$ of the \emph{entire prequantum groupoid} $(\TT_\omega, \blambda)$ that preserve the prequantum 1-form $\blambda$.
    The object being acted upon without central extension here is the groupoid's space of morphisms $\cY = \Paths(\X)/\!\sim_\omega$ endowed with the form $\blambda$.
    
    The classical prequantum bundle $\Y \to \X$
    (or, in our framework, the generalized prequantum bundle $\cY_x = \Mor(\TT_\omega, x, \ast) \to \X$, obtained by fixing a source point $x$)
    is \emph{not} the same object as the full groupoid $\TT_\omega$.
    The groupoid $\TT_\omega$ encompasses all such principal bundles $\cY_x$ for varying $x$,
    and the arrows of $\TT_\omega$ can be seen as isomorphisms between fibers of these bundles over different points.
    
    The central extensions observed in the classical setting \emph{will reappear} when one descends from the groupoid level to consider actions on structures derived from the groupoid over the base space $\X$.
    For example,
    when considering the lifting of symmetries to the total space of the principal bundle $\cY_x \to \X$ or,
    crucially,
    to the space of sections of an associated bundle over $\X$ (which would be the candidate quantum Hilbert space),
    the central extension structure related to the period group $\T_\omega$ will indeed typically emerge.
    
    Therefore,
    the result $\Aut(\TT_\omega, \blambda) \simeq \Diff(\X, \omega)$ does not contradict classical findings.
    Instead,
    it provides a potentially deeper insight: the groupoid $(\TT_\omega, \blambda)$ serves as a fundamental,
    extension-free prequantum object where the entire symmetry group $\Diff(\X, \omega)$ is faithfully represented.
    The "quantic anomaly" manifested by central extensions in the classical picture is then understood as a phenomenon that arises when one projects or realizes this fundamental structure and its symmetries within the context of bundles or representation spaces built over the base space $\X$.
    The essential quantic information,
    captured by the period group $\T_\omega$, resides fundamentally in the isotropy group of the groupoid,
    waiting to express itself as an extension upon descent to structures defined over $\X$. This perspective offers significant conceptual clarity,
    especially when dealing with singular spaces where classical bundle constructions are challenging.
    
    \smallskip
    \textbf{Note.} This phenomenon of symmetry breaking is reminiscent of passing from an equivariant two-point moment map $\psi(x,x')$ to a no longer equivariant one-point moment map $\mu(x) = \psi(x_0,x) + c$,
    which is the general solution of the equation $\psi(x,x') = \mu(x') - \mu(x)$.
    The one-point moment map is no longer equivariant under the linear coadjoint action but with the affine coadjoint action modified by the cocycle $\theta(g) = \psi(x_0,g(x_0)) + \Delta c$.
  \end{article}
  
  
  %%%%%%%%%%%%%%%%%%%%%%%%%%%%%%%%%%%%%%%%%%%%%%%%%%%%%%%%%%
  %% MARK: art: Future Directions
  %%%%%%%%%%%%%%%%%%%%%%%%%%%%%%%%%%%%%%%%%%%%%%%%%%%%%%%%%%
  
  \begin{article}
    \artlabel[Future Directions]
    \label{art:Future Directions}
    
    This paper has established the construction of the prequantum groupoid for connected and simply connected diffeological spaces.
    This constitutes the first step in a broader program to reformulate geometric quantization within the framework of diffeology.
    
    The extension of this construction to the general case (non-simply connected spaces),
    as well as the transition to the algebraic framework,
    are the subject of the sequel to this work:
    \emph{Geometric Quantization by Paths, Part II: The General Case}.
    In that work,
    we address:
    \begin{enumerate}
      \item \textbf{Homotopy and Periods:} We show how the fundamental group and the non-trivial homotopy of the space generate "surfacic periods."
      By incorporating these into the definition of the \emph{Total Group of Periods},
      we construct a prequantum groupoid for arbitrary connected spaces,
      resolving the obstructions associated with non-trivial topology.
      \item \textbf{The Prequantum Algebra:} We transition from the geometry of the groupoid to the analysis of the system by constructing the \emph{Convolution Algebra} of the groupoid.
      This algebra constitutes the intrinsic observable algebra of the system,
      retaining the full covariance of the classical symmetries.
    \end{enumerate}
    
    Beyond these immediate results,
    the groupoid framework opens several avenues for long-term research:
    
    $\ast$~\textbf{The Feynman Connection.}
    The path-based nature of our construction suggests that the prequantum groupoid is the "geometric crystallization" of the Feynman path integral.
    Formalizing the measure-theoretic link between the groupoid structure and the path integral amplitude remains a major open question.
    
    $\ast$~\textbf{Polarization and the Hilbert Space.}
    To recover standard quantum mechanics,
    one must select a specific representation of the prequantum algebra.
    This requires adapting the theory of \emph{polarizations} (Lagrangian or complex) to the groupoid setting,
    interpreting them as the choice of a "screen" on which the quantum geometry is projected.
    
    $\ast$~\textbf{The Maslov Correction.}
    In the classical bundle approach,
    the metaplectic correction is required to obtain the correct spectrum (e.g., for the harmonic oscillator).
    Investigating how this topological correction manifests intrinsically within the space of paths and the structure of the groupoid is essential.
    
    $\ast$~\textbf{Coadjoint Orbits.}
    Applying this construction to the coadjoint orbits of infinite-dimensional Lie groups (such as the Virasoro group or groups of diffeomorphisms) could provide new insights into their quantization,
    where standard manifold techniques often fail.
    
  \end{article}
  
  %%%%%%%%%%%%%%%%%%%%%%%%%%%%%%%%%%%%%%%%%%%%%%%%%%%%%%%%%%
  %% MARK: Bibliography
  %%%%%%%%%%%%%%%%%%%%%%%%%%%%%%%%%%%%%%%%%%%%%%%%%%%%%%%%%%
  
  \begin{thebibliography}{IKZ10} % Use 99 or a number larger than your total entries
    
    \bibitem[DI83]{DI83}
    Paul Donato and Patrick Iglesias.
    {\em Exemple de groupes diff\'erentiels : flots irrationnels sur le tore}.
    Preprint CPT-83/P.1524, Centre de Physique Th\'eorique, Marseille, July 1983. Published in {\em Comptes Rendus de l'Acad\'emie des Sciences}, 301(4), Paris, 1985.
    \texttt{http://math.huji.ac.il/~piz/documents/EDGDFISLT.pdf}
    
    \bibitem[GIZ23]{GIZ23}
    Serap G{\"u}rer and Patrick Iglesias-Zemmour.
    {\em Orbifolds as stratified diffeologies}.
    {Differential Geometry and its Applications}, 86:101969, 2023.
    
    \bibitem[GIZ25]{GIZ25}
    Serap Gürer and Patrick Iglesias-Zemmour,
    \emph{On Diffeology of Orbit Spaces},
    (Preprint available at \texttt{http://math.huji.ac.il//~piz/documents/ODOOS.pdf}).
    
    \bibitem[Hor24]{Hor24}
    Peter A. Horvathy.
    {\em Prequantisation from the path integral viewpoint}.
    Preprint arXiv:2402.17629, 2024.
    
    \bibitem[PIZ85]{PIZ85}
    Patrick Iglesias.
    {\em Fibr\'es diff\'eologiques et homotopie}.
    Th\`ese de doctorat d'\'etat, Universit\'e de Provence, Marseille, 1985.
    
    \bibitem[PIZ95]{PIZ95}
    Patrick Iglesias.
    {\em La trilogie du moment}.
    Ann. Inst. Fourier, 45(3):825--857, 1995.
    
    \bibitem[IKZ10]{IKZ10}
    Patrick Iglesias, Yael Karshon, and Moshe Zadka.
    {\em Orbifolds as diffeologies}.
    Transactions of the American Mathematical Society, Vol. 362, Number 6, Pages 2811--2831, Providence R.I., 2010.
    
    \bibitem[PIZ10]{PIZ10}
    Patrick Iglesias-Zemmour,
    \emph{The moment maps in diffeology},
    Mem. Amer. Math. Soc. \textbf{204} (2010), no. 962, viii+101 pp.
    
    \bibitem[PIZ13]{PIZ13}
    Patrick Iglesias-Zemmour.
    {\em Diffeology}.
    Mathematical Surveys and Monographs, 185, Am. Math. Soc., Providence, RI, 2013.
    (Revised version by Beijing World Publishing Corp., Beijing, 2022.)
    
    \bibitem[PIZ22]{PIZ22}
    Patrick Iglesias-Zemmour.
    {\em Diffeology}.
    Reprint by Beijing World Publishing Corporation, Ltd. Beijing Branch, 2022. ISBN 978-7-5192-9608-7.
    Accessible at \\
    \url{http://math.huji.ac.il/~piz/documents/Diffeology.pdf}
    
    \bibitem[PIZ15]{PIZ15}
    \bysame
    {\em Example of singular reduction in symplectic diffeology}.
    Proceedings of the American Mathematical Society, 144(3):1309--1324, 2016.
    
    \bibitem[IZP21]{IZP21}
    Patrick Iglesias-Zemmour and Elisa Prato.
    {\em Quasifolds, diffeology and noncommutative geometry}.
    J. Noncommut. Geom., 15(2):735--759, 2021.
    
    \bibitem[PIZ24]{PIZ24}
    Patrick Iglesias-Zemmour,
    \emph{{\v C}ech-de Rham bicomplex in diffeology},
    Israel J. Math. \textbf{259} (2024), no. 1, 239--276.
    
    \bibitem[PIZ25]{PIZ25}
    \bysame,
    \emph{The Boman Paradox --- When topology plus algebra do not define geometry},
    To appear in the Mathematical Intelligencer. Preprint 2025. Available at \url{http://math.huji.ac.il/~piz/documents/TBP.pdf}.
    
    \bibitem[MW74]{MW74}
    Jerrold E. Marsden and Alan Weinstein,
    \emph{Reduction of symplectic manifolds with symmetry},
    Rep. Mathematical Phys. \textbf{5} (1974), no. 1, 121--130.
    
    \bibitem[Sat56]{Sat56}
    Ichiro Satake,
    \emph{On a generalization of the notion of manifold},
    Proceedings of the National Academy of Sciences \textbf{42} (1956), 359--363.
    
    \bibitem[Sou70]{Sou70}
    Jean-Marie Souriau.
    {\em Structure des syst\`emes dynamiques}.
    Dunod, Paris, 1970.
    
    \bibitem[Sou75]{Sou75}
    \bysame %Jean-Marie Souriau.
    \newblock \emph{Construction explicite de l'indice de Maslov. Applications}.
    \newblock In: Group theoretical methods in physics (Fourth Internat. Colloq., Nijmegen, 1975), Lecture Notes in Phys., Vol. 50, Springer, Berlin, 1976, pp. 117--148.
    
    \bibitem[Sou04]{Sou04}
    Jean-Marie Souriau,
    \emph{C’est Quantique~? donc c’est géométrique},
    in: Dominique Flament (dir),
    \emph{Feuilletages -quantification géométrique : textes des journées d’étude des 16 et 17 octobre 2003, Paris, Fondation Maison des Sciences de l’Homme},
    Série Documents de travail (Équipe F2DS), 2004.
    
  \end{thebibliography}
  
  %%%%%%%%%%%%%%%%%%%%%%%%%%%%%%%%%%%%%%%%%%%%%%%%%%%%%%%%%%
  %%
  %% MARK: End Private comments
  %%
  %%%%%%%%%%%%%%%%%%%%%%%%%%%%%%%%%%%%%%%%%%%%%%%%%%%%%%%%%%
  
  
\end{document}
